% ============================================================
% ODE IVP Project — LaTeX Report Template
% Compile: pdflatex report.tex  (twice for the TOC)
% Structure follows the project brief: Background -> Model ->
% Methods -> Stability -> Numerical experiments -> Conclusions.
%
% This is the MERGED template: it combines the project-specific
% guidance for Project 1 (ODE IVPs) with the course-level report
% skeleton (algorithm environment, theorem environments, MMS
% convergence table, AI Transparency Log and ICS appendices).
% ============================================================

\documentclass[11pt,a4paper]{article}

\usepackage[utf8]{inputenc}
\usepackage[T1]{fontenc}
\usepackage{amsmath, amssymb, amsthm}
\usepackage{geometry}
\geometry{a4paper, margin=2.5cm}
\usepackage{graphicx}
\usepackage{booktabs}
\usepackage{listings}
\usepackage{xcolor}
\usepackage{hyperref}
\usepackage{enumitem}
\usepackage{float}
\usepackage{caption}
\usepackage{subcaption}
\usepackage{bm}
\usepackage{algorithm}
\usepackage{algpseudocode}
\usepackage{mdframed}

% --- colors ---
\definecolor{lecture}{RGB}{70,130,180}
\definecolor{seminar}{RGB}{205,92,92}
\definecolor{formbg}{RGB}{250,250,245}
\definecolor{formframe}{RGB}{180,170,140}

% --- code listing style (Python) ---
\lstset{
    language=Python,
    basicstyle=\ttfamily\scriptsize,
    keywordstyle=\color{blue}\bfseries,
    commentstyle=\color{green!40!black},
    stringstyle=\color{red!70!black},
    numbers=left,
    numberstyle=\tiny\color{gray},
    frame=single,
    breaklines=true,
    showstringspaces=false,
    tabsize=4,
    captionpos=b,
}

% --- theorem-like environments (course-level set) ---
\newtheorem{theorem}{Theorem}[section]
\newtheorem{lemma}[theorem]{Lemma}
\newtheorem{proposition}[theorem]{Proposition}
\newtheorem{corollary}[theorem]{Corollary}
\theoremstyle{definition}
\newtheorem{definition}[theorem]{Definition}
\newtheorem{assumption}[theorem]{Assumption}
\newtheorem{example}[theorem]{Example}
\theoremstyle{remark}
\newtheorem{remark}[theorem]{Remark}

% --- algorithm style ---
\algrenewcommand{\algorithmicrequire}{\textbf{Input:}}
\algrenewcommand{\algorithmicensure}{\textbf{Output:}}
\algrenewcommand{\algorithmiccomment}[1]{\hfill\textcolor{gray}{\# #1}}

% --- form box (for appendices) ---
\newmdenv[
    backgroundcolor=formbg,
    linecolor=formframe,
    linewidth=0.8pt,
    roundcorner=3pt,
    innermargin=5pt,
    outermargin=5pt,
    innertopmargin=8pt,
    innerbottommargin=8pt
]{formbox}

\title{\textbf{Numerical Solution of \texttt{<YOUR PROBLEM>}}\\[2mm]
\large Project 1 (ODE IVP) Report}
\author{
    Team \#\texttt{<N>} \\
    \small \texttt{<name1>, <id1> (Project Manager)} \quad
    \texttt{Meichen Ding, 2362344 (Mathematical Theory)} \\
    \small \texttt{<name3>, <id3> (Algorithm Implementation)} \quad
    \texttt{<name4>, <id4> (Visualization \& Report)} \\
    \small \texttt{<name5>, <id5> (Testing \& Validation)} \\
    \small Department of Applied Mathematics
}
\date{\today}

\begin{document}
\maketitle

\begin{abstract}
\noindent
[150--250 words answering: (1) what problem and why interesting; (2) what methods
you implemented; (3) key findings (convergence rates, stability observations);
(4) the main limitation or open question.]
\end{abstract}

\noindent\textbf{Keywords:} [e.g., Runge--Kutta methods, stiffness, von Neumann
stability, absolute stability, adaptive stepping]

\tableofcontents
\newpage

% ============================================================
\section{Introduction and Motivation}
% ============================================================

\subsection{Problem Background}
[Describe the physical, biological, or financial context of your chosen ODE
system. Why is this problem worth solving? Cite 2--3 references.]

\subsection{Mathematical Model}
[State the ODE system, initial conditions, and all parameter values. If the
system is stiff, say so and give the stiffness ratio.]

\begin{equation}
    \frac{dy}{dt} = f(t, y), \qquad y(t_0) = y_0.
    \label{eq:ivp}
\end{equation}

For the Robertson problem (example):
\begin{equation}
\begin{aligned}
    \dot{y}_1 &= -0.04\, y_1 + 10^4\, y_2 y_3, \\
    \dot{y}_2 &= \;\;0.04\, y_1 - 10^4\, y_2 y_3 - 3\times 10^7\, y_2^2, \\
    \dot{y}_3 &= \;\;3\times 10^7\, y_2^2,
\end{aligned}
\qquad y(0) = (1, 0, 0).
\label{eq:robertson}
\end{equation}

\subsection{Overview of the Report}
[A roadmap: ``In Section 2 we derive the discrete scheme and analyse stability.
In Section 3 we describe the implementation. Section 4 presents numerical
experiments. Section 5 concludes.'']

\newpage
% ============================================================
\section{Discretisation and Analysis}
% ============================================================

\subsection{Discrete Schemes}
[Derive your three methods step by step: Explicit Euler, RK4 (show the Butcher
tableau), and your implicit method with Newton. Number every equation.]

\begin{example}[ODE IVP: RK4 derivation]
For $\dot{y} = f(t,y)$, the classical Runge--Kutta method of order 4 is
\begin{align}
    k_1 &= f(t_n, y_n), \\
    k_2 &= f\!\left(t_n + \frac{h}{2}, y_n + \frac{h}{2}k_1\right), \\
    k_3 &= f\!\left(t_n + \frac{h}{2}, y_n + \frac{h}{2}k_2\right), \\
    k_4 &= f(t_n + h, y_n + hk_3), \\
    y_{n+1} &= y_n + \frac{h}{6}(k_1 + 2k_2 + 2k_3 + k_4).
\end{align}
\end{example}

\subsection{Consistency and Truncation Error}
[Show the Taylor expansion. State the order of accuracy of each method and under
what assumptions it holds.]

\subsection{Stability Analysis}
[The most critical section. Do not quote a textbook result --- apply the analysis
to \emph{your} problem. Derive the stability function $R(z)$ of each method, plot
the absolute-stability region, and connect it to the eigenvalues of your
linearised system.]

\begin{equation}
    R_{\text{Euler}}(z) = 1 + z, \qquad
    R_{\text{RK4}}(z) = 1 + z + \frac{z^2}{2} + \frac{z^3}{6} + \frac{z^4}{24},
    \qquad
    R_{\text{impl.Euler}}(z) = \frac{1}{1 - z}.
\end{equation}

\subsection{Convergence}
[Link consistency + stability to convergence (Lax equivalence or an energy
argument). State the theorem and verify your problem satisfies its hypotheses.]

\newpage
% ============================================================
\section{Implementation}
% ============================================================

\subsection{Algorithm}
[Present the core algorithm in pseudocode using the \texttt{algorithmic}
environment.]

\begin{algorithm}[H]
\caption{[Your method name]}
\label{alg:main}
\begin{algorithmic}[1]
\Require [initial data $y_0$, step size $h$, final time $T$]
\Ensure [approximate solution at $T$]
\State [Step 1]
\For{$n = 0, 1, \ldots, N-1$}
    \State [Loop body]
    \If{[condition]}
        \State [Action]
    \EndIf
\EndFor
\State \Return [result]
\end{algorithmic}
\end{algorithm}

\subsection{Code Structure}
[Briefly describe the organisation of your codebase. A table of files and their
purposes is helpful.]

\begin{table}[H]
\centering
\begin{tabular}{ll}
\toprule
\textbf{File} & \textbf{Purpose} \\
\midrule
\texttt{methods.py} & Core time-stepping: Euler, RK4, implicit Euler with Newton \\
\texttt{model.py} & The right-hand side $f(t,y)$ and its Jacobian \\
\texttt{adaptive.py} & Adaptive step-size controller \\
\texttt{run\_all.py} & Reproduces all report figures \\
\texttt{validation\_tests.py} & PASS/FAIL test suite \\
\bottomrule
\end{tabular}
\end{table}

\subsection{Key Design Decisions}
[Discuss non-obvious choices: why this implicit method, how Newton's initial
guess is chosen, how the adaptive controller estimates local error.]

\newpage
% ============================================================
\section{Numerical Experiments}
% ============================================================

\subsection{Verification: Method of Manufactured Solutions or Exact Solutions}
[Present a test case with a known solution. Show the error table and the
convergence plot.]

\begin{table}[H]
\centering
\begin{tabular}{cccc}
\toprule
\textbf{Step size $h$} & \textbf{Error} & \textbf{Ratio} & \textbf{Order} \\
\midrule
$0.1$    & [value] & ---  & ---  \\
$0.05$   & [value] & [ratio] & [order] \\
$0.025$  & [value] & [ratio] & [order] \\
\bottomrule
\end{tabular}
\caption{Convergence study for [test case].}
\end{table}

\subsection{Stability Experiments}
[Show that your explicit methods blow up beyond their stability limit while the
implicit method does not. Compare with the predicted threshold.]

\subsection{Adaptivity}
[Show the step size actually changes, the local error stays near the tolerance,
and the $h(t)$ trajectory.]

\subsection{Main Results and Comparison of Methods}
[Present the experiments that answer the project's central question. Use
high-quality figures with interpretive captions.]

\begin{figure}[H]
\centering
\fbox{\parbox{0.8\textwidth}{\centering \vspace{3em} [Insert your figure here] \vspace{3em}}}
\caption{[Caption that interprets, not merely describes. E.g., ``The log-log plot
confirms second-order convergence (slope 1.98) on the smooth test case, and shows
the round-off floor at $h \approx 10^{-6}$.'']}
\label{fig:convergence}
\end{figure}

\subsection{Limitations and Robustness}
[Discuss parameter regimes where your method struggles. This is not a weakness ---
it is scientific honesty.]

\newpage
% ============================================================
\section{Conclusions}
% ============================================================

[Summarise in 3--4 bullet points:]
\begin{itemize}[leftmargin=2em]
\item [What you achieved.]
\item [The most surprising or instructive finding.]
\item [The main limitation you identified.]
\item [A concrete direction for future work.]
\end{itemize}

\newpage
% ============================================================
\section*{References}
% ============================================================

[Use BibTeX or list references manually. Include at least: (1) the textbook from
which you learned the method; (2) a paper applying the method to a problem
similar to yours; (3) a software reference (NumPy, SciPy, scikit-fem).]

\begin{thebibliography}{9}
\bibitem{ref1} Author, A. \emph{Title}. Journal, year.
\bibitem{ref2} Author, B. \emph{Book title}. Publisher, year.
\end{thebibliography}

\newpage
% ============================================================
\appendix
\section{AI Transparency Log}
% ============================================================

[Append the completed team AI Transparency Log here. See
\texttt{06\_Templates/ICS} for the full form: tools used, significant claim audits,
what you changed and why, and the intellectual ownership declarations.]

\section{Individual Contribution Statements}
% ============================================================

Each team member must complete the following (repeat for all 5 members):

\begin{formbox}
\textbf{Name:} Meichen Ding \quad (\textbf{Student ID:} 2362344) \hfill \textbf{Role:} Mathematical Theory Lead

\vspace{0.5em}
\textbf{1. My specific contributions to this project:}

I authored the analytical consistency, adaptivity, and convergence foundations in Section~2 (Sections~2.2 and~2.4) now in \path{Project1/report_v3/sections/section2.tex}; an earlier draft is preserved as \path{Project1/report_v1/Section2Draft_v5.tex}. Specifically, I derived the local truncation errors and consistency orders for Explicit/Implicit Euler and RK4 (Eq.~37, 43, 45--52 in Section2Draft\_v5) under Picard--Lindel{\"o}f Lipschitz continuity and small-step contraction $hL < 1$ (Eq.~42). From first principles, I derived the step-doubling Richardson extrapolation error estimator, proving that the exact continuous state cancels out to yield the unnormalised defect denominator $2^p - 1$ (Eq.~57--58) and the scale-invariant step-size update controller with optimal exponent $-1/(p+1)$ (Eq.~62). Furthermore, I formulated the discrete error recurrence (Eq.~104--105) and proved the continuous-type discrete Gr{\"o}nwall global upper bound $\|e_n\| \le \frac{e^{\Lambda(T-t_0)}-1}{\Lambda}\tau_{\max}(h)$ (Eq.~109), establishing why local defects of order $\mathcal{O}(h^{p+1})$ accumulate over $N=\mathcal{O}(h^{-1})$ steps into $\mathcal{O}(h^p)$ global convergence (Eq.~110--111), while proving the zero-stability degeneracy ($\rho(\zeta)=\zeta-1=0$) of single-step methods and providing the theoretical derivation for empirical log-log slope regression (Eq.~114--115).

\vspace{0.5em}
\textbf{2. The hardest conceptual obstacle I encountered:}

The hardest mathematical challenge was reconciling why accepting the fine sub-step solution $y_f$ in our step-doubling adaptive controller did not elevate the asymptotic global convergence order from $p$ to $p+1$ in empirical log-log regressions, despite the controller utilizing a local extrapolation error formula. I overcame this by writing out the full Taylor remainder expansions for both the coarse step $y_c$ and consecutive half-steps $y_f$, which mathematically proved that combining two half-steps merely attenuates the leading principal error coefficient by a factor of $2^{-p}$ without satisfying the underlying Butcher algebraic tree conditions required for order $p+1$, thereby clarifying that local error control must strictly follow the exponent $-1/(p+1)$.

\vspace{0.5em}
\textbf{3. One piece of AI-generated output my team used, and what I changed:}

Generative AI was consulted to revise the Implicit Euler local-error argument and adaptive stepping. When querying AI for the adaptive step-size update rule, the raw output proposed a textbook formula $h_{\text{new}} = h \cdot E^{-1/p}$. I audited this claim and identified that it conflated accumulated global discretization error $\mathcal{O}(h^p)$ with single-step local truncation defect $\mathcal{O}(h^{p+1})$. I corrected the resizing exponent to $-1/(p+1)$ in Section~2.2.4 (Eq.~62), derived the rigorous step-doubling cancellation denominator $2^p - 1$ ($1$ for Euler, $15$ for RK4), and introduced the safety margin $S=0.9$ alongside limiter bounds $[r_{\min}, r_{\max}]$ to prevent step-chattering during the initial transient.

\vspace{0.5em}
\textbf{4. One thing I still do not fully understand:}

I do not yet fully understand how to construct a rigorous non-autonomous matrix logarithmic norm (one-sided Lipschitz dissipative constant) for the nonlinear polar-factor gradient flow $\dot{X} = X(I_n - X^TX)$ over the entire finite-time trajectory $t \in [0, 2]$, which would provide a unified global nonlinear contraction bound to theoretically explain why the discrete Lyapunov residual exhibits localized pre-asymptotic spikes during the fast transient even when the continuous Lyapunov functional $\Phi(X)$ decreases monotonically.

\vspace{0.8em}
\textbf{Signature:} \underline{\hspace{4.5cm}} \qquad \textbf{Date:} \underline{\hspace{2.5cm}}
\end{formbox}Meichen

\section{Supplementary Figures and Data}
% ============================================================

[Place any additional figures, tables, or code listings that would interrupt the
main narrative but may be useful to a curious reader.]

\end{document}
